\chapter{Introduction}
\section{Project Brief}
\subsection{Motivation}
The pharmaceutical industry faces high costs and long development cycles in drug discovery, partly due to the challenge of identifying whether a novel compound can bind to a specific protein target.
Machine learning (ML) provides a promising pathway for automating and accelerating this process by learning patterns from large bioactivity datasets.
\subsection{Objective}
The aim of this project is to develop a predictive machine learning system that estimates the likelihood of binding between a drug-like molecule and a protein target during the early stages of discovery.
This system is intended to reduce the sample size taken forward to wet-lab validation to hopefully shorten the development cycles the pharmaceutical industry faces, by removing confident non-binders.
\subsection{Terminology}
Throughout this report the terms \textit{drug}, \textit{compound}, and \textit{small molecule} are used interchangeably to refer to a ligand being evaluated against a protein target.
A \textit{drug} is a compound with clinical approval, but in the virtual-screening context every candidate is treated uniformly by the model and no distinction is required.
Similarly, \textit{protein}, \textit{target}, and \textit{protein target} refer to the same biological entity against which binding is predicted; \textit{target} emphasises the therapeutic intent, while \textit{protein} is the structural description, and the two are used interchangeably.
\subsection{Computational Approaches in DTI}
Drug-target interaction was investigated and tested through wet-lab experiments until the first \textit{in silico} approach, molecular docking against 3D structures, was introduced in 1982~\citep{Ru2025InSilicoDTI}.
A later breakthrough came with the adoption of machine learning models from the early 2000s onwards, which can autonomously recognise patterns and relationships in bioactivity data without requiring explicit physical modelling.
A significant step forward followed in 2018 \citep{Ozturk2018}, when deep learning architectures were shown to learn effective representations directly from drug and protein sequence data.
\subsection{Process}
\subsubsection{Data Acquisition}
Relevant bioactivity data was extracted and processed from BindingDB, a public database containing hundreds of thousands of drug-target binding measurements~\citep{Liu2025}.
\subsubsection{Feature Engineering}
Compounds were converted to molecular fingerprints using RDKit, and protein targets were encoded as numerical vectors using ProtBERT.
These pipelines differ because the two modalities encode biological information in fundamentally different forms, requiring domain-specific feature engineering.
\subsubsection{Model Development, Training \& Evaluation}
A neural network was developed and trained on the processed dataset to learn the underlying relationships that distinguish binders from non-binders.
The model was evaluated using AUC-ROC, PR-AUC, F1, precision, recall, and accuracy.
Based on these metrics, the network was iteratively refined to reduce error and improve discriminative performance.
\subsubsection{Develop User Interface}
The website was developed using the React framework for its modular architecture, rapid development workflow, and live rendering~\citep{aggarwal2018react}.
The interface was designed to be easy to navigate and to present only the information relevant to the workflow.
It displays each prediction result, a history of previous comparisons, and an option to compare multiple compounds against a single protein to mimic wet-lab virtual screening.
\subsubsection{Tech Stack}
\begin{table}[H]
    \centering
    \begin{tabular}{@{}lp{10cm}@{}}
    \toprule
    \textbf{Component} & \textbf{Technology / Library} \\
    \midrule
    Dataset & BindingDB  \\
    \addlinespace
    \midrule
    Programming Language & Python, JavaScript\\
    \addlinespace
    \midrule
    Data Processing & Pandas, NumPy, Scikit-Learn \\
    \addlinespace
    \midrule
    Drug Molecule Feature Engineering & RDKit (molecular fingerprints, SMILES parsing, descriptor extraction) \\
    \addlinespace
    \midrule
    Protein Feature Engineering & ProtBERT (protein embeddings using transformer-based model) \\
    \addlinespace
    \midrule
    Neural Network Framework & PyTorch (model definition, training, and evaluation) \\
    \addlinespace
    \midrule
    Model Evaluation & Scikit-learn \\
    \addlinespace
    \midrule
    User Interface & React (frontend framework for model interaction and visualization) \\
    \addlinespace
    \midrule
    Backend API & FastAPI (model serving and communication with frontend) \\
    \addlinespace
    \midrule
    Version Control & Git / GitHub \\
    \addlinespace
    \bottomrule
    \end{tabular}
    \caption{Technologies and tools used in the project.}
    \label{tab:tech_stack}
\end{table}

\section{Feature List}
This section discusses the features planned for the project and summarises the role of each within the overall workflow.

\subsection{Required}
\subsubsection{Compare a single drug molecule to a single protein target}
The most basic functionality of the system: the user searches the database for a single compound and tests it against a single protein target.
\subsubsection{Compare multiple drug molecules to a single protein target}
The user can multi-select compounds, or filter them by physicochemical properties and known activity, and screen the resulting set against a chosen target.
This mirrors the triage process in wet-lab virtual screening.
\subsubsection{Displaying previous tests}
A prediction history reduces redundant computation by allowing the user to revisit earlier results rather than re-running the network for the same drug--target pair.
\subsection{Nice-to-Have}
\subsubsection{Produce indirect links into how they bind}
Once a prediction has been returned, a brief explanation of the features driving the outcome would help users reason about why the network produced that score.
\subsubsection{Displaying molecule and Protein information}
Retrieving contextual metadata from BindingDB allows the user to inspect additional properties of the selected drug-protein pair alongside the prediction.

\section{Scope \& Limitations}
\subsection{Project Scope}
The project was broken into milestones so that each reached state represented a self-contained, submittable piece of work forming the foundation for the next:
\subsubsection{Predicting drug-target interactions}
A NumPy-based model capable of classifying drug-target pairs as binding or non-binding.
\subsubsection{Using neural networks for prediction}
A neural network architecture trained to learn the relationships between molecule and protein features and produce consistent predictions.
\subsubsection{Including both chemical compound and protein sequence features}
Molecular descriptors or fingerprints for compounds, combined with sequence-based or embedding-based features for protein targets.
\subsubsection{Implementing a web-based UI for result visualization and comparison}
A React-based web application that visualises prediction results, supports drug-target comparisons, and provides basic interpretability.

\subsection{Project Limitations}
The following factors constrain the scope of the project and the conclusions that can be drawn from it.
\subsubsection{Data limitations}
The project is restricted to only using publicly available datasets (e.g.\ BindingDB, ChEMBL), which may contain missing or biased entries introduced during curation.
This can affect the model's predictive performance if the dataset isn't complete.
\subsubsection{Model limitations}
Fully connected neural networks may struggle with rare interactions or out-of-distribution compounds, and they do not encode the 3D structure of each molecule.
A graph neural network (GNN) could be used instead, but that would introduce additional computational and data requirements beyond the resources available for this project.
\subsubsection{Computational constraints}
The model was trained on a personal laptop, so the computational limits are not representative of what industrial pharmaceutical research groups have, widening the gap between this system and industrial systems.
\subsubsection{UI limitations}
The interface is a prototype intended for small-team use rather than production deployment.
\subsection{Assumptions}
The project assumes that users has sufficient chemistry knowledge to interpret the input fields and results.
It also assumes that the data within BindingDB is accurate and up to date.

\section{Workflow}
This section presents a visual representation of the project workflow from data acquisition through to a deployed user interface.
\subsection{System Overview}
The system predicts potential interactions between compounds and protein targets using a trained neural network.
It processes input data, extracts numerical features, performs inference, and displays the results through a web interface for comparison and interpretation.
\begin{table}[h!]
    \centering
    \begin{tabular}{@{}lp{10cm}@{}}
    \toprule
    \textbf{Component} & \textbf{Description} \\
    \midrule
    Data Collection \& Preprocessing & Collecting compound and protein data from public databases such as BindingDB, and filtering it based on the required parameters. \\
    \addlinespace
    \midrule
    Feature Engineering & Converting molecular structures and protein sequences into numerical feature representations such as fingerprints \& embeddings \\
    \addlinespace
    \midrule
    Model Training and Evaluation & Training a neural network on the processed dataset to pick up on relationships and evaluating its performance using metrics such as AUC-ROC and F1 score. \\
    \addlinespace
    \midrule
    User Interface & Displays results and prediction history, and allows the user to customise comparisons. Built with React for responsiveness and ease of use. \\
    \bottomrule
    \end{tabular}
    \caption{System components and their descriptions.}
    \label{tab:system_components}
\end{table}
\subsection{Workflow Diagram}

\begin{figure}[H]
    \hspace{0.8725cm}
    \begin{tikzpicture}[
        % Define styles for the nodes
        main_process/.style={rectangle, minimum width=4cm, minimum height=1cm, text centered, draw=black, fill=blue!30, line width=0.5pt, font=\bfseries},
        sub_process/.style={rectangle, minimum width=3.5cm, minimum height=0.8cm, text centered, draw=black, fill=green!30, line width=0.5pt},
        arrow/.style={-Latex, very thick, line width=1pt},
        feedback/.style={-Latex, thick, dashed, line width=1pt, color=red!80!black} % For the testing loop
    ]

    % 1. Data Collection Node
    \node (DC) [main_process, align=center] {Data Collection};

    % 2. Feature Engineering Node (Parent)
    \node (FE_Parent) [main_process, below=1.5cm of DC] {Feature Engineering};

    % 3. Spiked-off Sub-nodes
    \node (MF) [sub_process, below=1.25cm of FE_Parent, xshift=-2.75cm] {Molecular Fingerprints \\ (RDKit)};
    \node (NV) [sub_process, below=1.25cm of FE_Parent, xshift=2.75cm] {Numerical Vectors \\ (ProtBERT)};

    % 4. Model Training & Evaluation Node
    % Positioning it below the midpoint of the sub-nodes
    \node (MTE) [main_process, below=3.5cm of FE_Parent] {Model Training \\ \& Evaluation};

    % 5. Web UI Deployment Node
    \node (WEB) [main_process, below=1.5cm of MTE] {Web UI Deployment};


    % --- Draw Arrows ---

    % Step 1 -> Step 2
    \draw [arrow] (DC) -- (FE_Parent);

    % Step 2 Spiking Off
    \draw [arrow] (FE_Parent) -- (MF);
    \draw [arrow] (FE_Parent) -- (NV);

    % Sub-nodes return to Step 3
    \draw [arrow] (MF) -- (MTE);
    \draw [arrow] (NV) -- (MTE);



    % self loop:
    \draw [feedback] (MTE.south west)
        -- ++(-2cm, 0)
        |- (MTE.north west)
        node[midway, left=0cm, yshift=-0.5cm] {Testing};
    



    % Step 3 -> Step 4
    \draw [arrow] (MTE) -- (WEB);

    \end{tikzpicture}
    \caption{Top-down project workflow illustrating parallel feature engineering and iterative model refinement.}
    \label{fig:workflow_diagram_v2}
\end{figure}

\subsection{Project Timeline}
A Gantt chart was used to determine major milestones and track progress across the development cycle.
\begin{figure}[H]
    \centering
    \includegraphics[width=\textwidth, trim = 1cm 1cm 483.75cm 1cm, clip]{IC-12-Month-Project-Gantt-Chart.pdf}
    \caption{Project Gantt chart showing major milestones and development phases.}
    \label{fig:gantt_chart}
\end{figure}